\documentclass[11pt]{article}

\usepackage[utf8]{inputenc}
\usepackage[T1]{fontenc}
\usepackage{lmodern}
\usepackage{geometry}
\usepackage{amsmath,amssymb,amsfonts,amsthm,mathtools}
\usepackage{bm}
\usepackage{enumitem}
\usepackage{microtype}
\usepackage{hyperref}
\usepackage{titlesec}
\usepackage{booktabs}
\usepackage{array}
\usepackage{setspace}
\usepackage{mathrsfs}
\usepackage{physics}

\geometry{margin=1in}
\hypersetup{colorlinks=true, linkcolor=black, urlcolor=blue, citecolor=black}
\setlist[enumerate]{topsep=0.4em,itemsep=0.35em}
\setlist[itemize]{topsep=0.3em,itemsep=0.25em}
\titleformat{\section}{\large\bfseries}{\thesection.}{0.5em}{}
\titleformat{\subsection}{\normalsize\bfseries}{\thesubsection.}{0.5em}{}
\setstretch{1.06}

\newcommand{\Aether}{\AE{}ther}
\newcommand{\Aflow}{\AE{}ther-flow}
\newcommand{\TheoryName}{\Aflow{} Interpretation of Relativity}
\newcommand{\TheoryProgram}{\Aether{} / \Aflow{} framework}
\newcommand{\CompletedMetric}{g^{\sharp}}

\newcommand{\ReducedStructure}{\mathfrak S^{\mathrm{disp,resp,red}}}
\newcommand{\MinimalReducedStructure}{\mathfrak S^{\mathrm{disp,resp,red,min}}}

\newtheorem{definition}{Definition}
\newtheorem{proposition}{Proposition}
\newtheorem{remark}{Remark}
\newtheorem{corollary}{Corollary}
\newtheorem{theorem}{Theorem}

\title{\textbf{The \TheoryName{}}\\[0.4em]
\large Primitive-Reservoir Observer-Reduction\\
\large Section-Centered Hidden-Response\\
\large Reduced Projection-Response Substrate Structure\\
\large Necessity / Uniqueness Next Benchmark Criterion}
\author{}
\date{}

\begin{document}

\maketitle

\begin{abstract}
The current active-primary theorem now proves that, once the fixed full canonical current-body image core is held fixed, every benchmark-faithful reduced projection-response substrate structure compatible with that core induces the same visible quadratic response operator on the fixed section-centered displacement fibers and differs from the canonical minimal representative \(\MinimalReducedStructure\) only by inert kernel directions. The reduced projection-response substrate structure is therefore benchmark-facingly unavoidable up to kernel-stable equivalence.

This note records the routing consequence of that gain. The next honest benchmark-facing move must now derive or justify the reduced projection-response substrate structure \(\ReducedStructure\) itself, equivalently its canonical minimal visible-response representative \(\MinimalReducedStructure\), from still more primitive explicit substrate structure while preserving the same benchmark one-metric package, the same ordinary matter route, and the same claim boundary. Another manuscript that merely adjoins inert kernel directions, renames the same visible-response representative, restores the larger substrate-body presentation, or replays the larger projection-operator or displacement-response presentations does not count as the clean next step. The one operative metric remains exactly \(\CompletedMetric\), there is no second operative metric, the low-energy matter law is not enlarged, and no stronger promotion language is licensed.
\end{abstract}

\section{Introduction}

The active derivational program of \TheoryName{} is no longer merely at full-image-core forcing scope. The preceding theorem already showed that once the reduced projection-response substrate structure is fixed, the full canonical current-body image core follows canonically and uniquely. The new reduced-structure necessity / uniqueness theorem now sharpens that gain one honest level further: once the fixed current image core is held fixed, every compatible reduced projection-response substrate structure collapses to the same visible-response representative, and the only remaining carrier-level freedom is inert kernel extension.

The present note records the routing consequence of reading those two theorems together. It does not reopen the already-screened larger substrate-body presentation, the older projection-operator core, the full displacement-response substrate law, or the now-spent option of changing the reduced structure merely by adjoining or removing compatible inert kernel directions. It records only which kind of next manuscript would now count as an honest benchmark-facing gain below the reduced-structure necessity / uniqueness frontier.

\paragraph{Framework and claim boundary.}
\TheoryProgram{} denotes the broader \Aether{} / \Aflow{} framework built on the claims that the \Aether{} is the underlying four-dimensional substrate of reality, the \Aflow{} is its intrinsic ordered motion, observed three-dimensional space is the local experiential slice of that deeper substrate, S-time is the experienced order of change arising from matter, light, and the \Aflow{}, and observed expansion is the three-dimensional appearance of deeper four-dimensional ordered motion. Within that framework, \TheoryName{} denotes the exact-closure theory in which the effective gravitational dynamics are adopted to be exactly Einsteinian with universal matter coupling. In the active sequence, \emph{adoption} means use of that established relativistic dynamics without claiming substrate derivation, while \emph{derivation} is reserved for a first-principles recovery from explicit substrate variables.

The present note stays strictly inside that boundary. It records only which kind of next manuscript would now count as an honest benchmark-facing gain below the reduced-projection-response-structure necessity / uniqueness frontier.

\section{Fixed Inputs}

\begin{definition}[Reduced-structure necessity / uniqueness frontier inputs]
The criterion recorded here uses the following fixed inputs.
\begin{enumerate}
    \item The benchmark exact-closure package, which fixes the public one-metric GR target of the repository.
    \item The benchmark gatekeeping layer, including the recorded safeguard that blocks same-output deeper-origin relays from counting as new front-line progress once the live benchmark-relevant datum is unchanged.
    \item The theorem proving that, once a reduced projection-response substrate structure \(\ReducedStructure\) is fixed, the full canonical current-body image core is canonically derived and uniquely forced.
    \item The later theorem proving that, once the fixed full canonical current-body image core is held fixed, every benchmark-faithful compatible reduced projection-response substrate structure collapses to the same canonical minimal representative \(\MinimalReducedStructure\) up to inert kernel extension.
\end{enumerate}
\end{definition}

\begin{remark}
The present note does not replace those manuscripts. It records the routing consequence of reading them together under the active safeguard against recursive depth drift.
\end{remark}

\section{Why the Live Burden Now Sits Below Compatible Carrier Restatement}

\begin{proposition}[The live burden is renewed below inert-kernel restatement]
At the current stage of the primitive-reservoir observer-reduction line, the reduced-structure necessity / uniqueness theorem renews the live benchmark-facing burden below any mere compatible carrier restatement of the same reduced projection-response substrate structure \(\ReducedStructure\) or of the same canonical minimal representative \(\MinimalReducedStructure\).
\end{proposition}

\begin{proof}
That theorem changes benchmark-facing status because it proves that once the fixed current image core is held fixed, every compatible reduced projection-response substrate structure carries the same visible-response representative and differs only by inert kernel directions. Once that uniqueness statement is in hand, another manuscript that merely adjoins kernel directions, removes them again, or reparameterizes the same visible-response representative no longer removes any remaining benchmark-facing freedom. The live burden therefore no longer sits on another compatible carrier presentation of the same reduced structure. It sits on deriving that reduced structure itself, equivalently its canonical minimal representative, from still more primitive explicit substrate structure.
\end{proof}

\begin{definition}[Benchmark-neutral inert-kernel extension relay]
A \emph{benchmark-neutral inert-kernel extension relay below the current frontier} is a manuscript that preserves the same benchmark one-metric output, the same ordinary matter route, the same reduced projection-response substrate structure \(\ReducedStructure\), the same canonical minimal representative \(\MinimalReducedStructure\), and the same full canonical current-body image core while merely adjoining or removing compatible inert kernel directions, renaming the same visible-response representative, restoring the already-screened larger substrate-body presentation, or re-expanding back to the older projection-operator core, the full displacement-response substrate law, or the larger pullback-body-operator, pullback-operator, pullback-quadratic, hidden-response, residual-normal-form, or pre-proto-section presentations without a new derivation of the reduced structure itself.
\end{definition}

\section{The Current Post-Necessity Criterion}

\begin{definition}[Admissible next benchmark-facing manuscript below the current frontier]
Below the current reduced-projection-response-structure necessity / uniqueness frontier, a manuscript counts as an admissible next benchmark-facing move only if it does at least one of the following:
\begin{enumerate}
    \item derives or justifies the reduced projection-response substrate structure \(\ReducedStructure\) itself from still more primitive explicit substrate structure;
    \item equivalently derives the canonical minimal visible-response representative \(\MinimalReducedStructure\) from still more primitive explicit substrate structure in a structure-faithful way;
    \item does one of the previous two without reopening the larger substrate-body presentation, the older projection-operator core, the full displacement-response substrate law, or another compatible inert-kernel extension as if any of those were again the live burden.
\end{enumerate}
\end{definition}

\begin{theorem}[Next honest benchmark-facing task below the current frontier]
The next honest benchmark-facing manuscript below the current reduced-projection-response-structure necessity / uniqueness frontier must derive or justify the reduced projection-response substrate structure \(\ReducedStructure\) itself, equivalently its canonical minimal representative \(\MinimalReducedStructure\), from still more primitive explicit substrate structure. A manuscript that only repeats the same reduced structure through another compatible carrier presentation, merely adjoins inert kernel directions, renames the same visible-response representative, restores the larger substrate-body presentation, or re-expands the discussion back to the older projection-operator core or the full displacement-response substrate law does not satisfy the criterion.
\end{theorem}

\begin{proof}
By Proposition 1, the live burden already lies below any mere compatible carrier restatement of the same reduced structure. Therefore a subsequent manuscript must improve on the already-forced reduced-structure datum itself rather than merely accompany it through another presentation. A deeper-origin manuscript can do so only if it changes benchmark-facing status by more than carrier-level relabeling or inert kernel extension. The only honest remaining task is therefore to derive the reduced structure itself, equivalently its canonical minimal representative, from still more primitive explicit substrate structure.
\end{proof}

\section{What the Next Move Should Not Be}

\begin{proposition}[Screened next moves below the current frontier]
Below the current reduced-projection-response-structure necessity / uniqueness frontier, none of the following is the default next benchmark-facing move:
\begin{enumerate}
    \item another same-output deeper-origin relay that merely preserves the current reduced projection-response substrate structure \(\ReducedStructure\) or merely presents the same canonical minimal representative \(\MinimalReducedStructure\) through a different compatible carrier and inert kernel extension and therefore reproduces the same full canonical current-body image core together with the same downstream benchmark-facing consequences;
    \item another restoration of the full substrate-body presentation as if the already-screened larger body freedom were again independently live;
    \item another replay of the older projection-operator core, the full displacement-response substrate law, or the larger pullback-body-operator, pullback-operator, pullback-quadratic, hidden-response, residual-normal-form, or pre-proto-section presentations as if those already-spent larger presentations were again the honest current burden;
    \item a manuscript that introduces a second operative metric, enlarges the ordinary matter law, or uses stronger promotion language.
\end{enumerate}
\end{proposition}

\begin{proof}
Item (1) fails the preceding criterion because it leaves the renewed live burden unchanged. Item (2) likewise fails because it restores a larger presentation whose relevant freedom is already screened out. Item (3) fails because it re-expands to larger already-spent presentations rather than deriving the reduced structure itself. Item (4) violates the fixed claim boundary of the benchmark package and therefore cannot count as a disciplined next step on the active line.
\end{proof}

\begin{remark}
This screening statement is not a denial that those deeper or alternative manuscripts may matter strategically. It is a routing statement: they are not the default way to spend the next benchmark-facing move below the current frontier unless they do the benchmark-facing work named above.
\end{remark}

\section{Conclusion}

The positive result of this note is routing discipline below the reduced-projection-response-structure necessity / uniqueness frontier. The present theorem remains the live benchmark-facing gain because it already proves that, once the fixed current image core is held fixed, every compatible reduced projection-response substrate structure collapses to the same canonical minimal representative \(\MinimalReducedStructure\) up to inert kernel extension. The earlier full-image-core forcing theorem remains preserved handoff history because it first made the reduced structure itself the live burden.

The scope remains narrow and honest. The benchmark package is unchanged: the one operative metric is still exactly \(\CompletedMetric\), the ordinary matter route is unchanged, there is no second operative metric, and the low-energy matter law is not enlarged. Nothing here upgrades adoption to derivation or licenses stronger promotion language.

The next open burden is therefore clear. The next benchmark-facing manuscript should derive or justify the reduced projection-response substrate structure \(\ReducedStructure\) itself from still more primitive explicit substrate structure, equivalently derive its canonical minimal visible-response representative \(\MinimalReducedStructure\). Another same-output relay that merely varies compatible kernel directions, another restoration of the full substrate-body presentation, or another replay of the older projection-operator core or the full displacement-response substrate law is not the clean next manuscript target at current scope.

\end{document}
